\begin{proposition}[independence of the equivariant lift]\label{push:src:C.3.6}
Let \(X=X_c\) carry the effective semifree circle action of \ref{push:src:C.3.1},
with fixed surface \(F\cong S^2\). If two smooth equivariant
diffeomorphisms of \(X\) induce the same map of its orbit space,
they are smoothly isotopic. More generally, a smooth isotopy of
orbit maps which is supported in a compact subset of the free
orbit stratum lifts to an equivariant isotopy. Thus, whenever the
annular orbit twist has been constructed with a smooth equivariant
lift, its class in \(\Theta_7\) is independent of that lift.
This assertion does not construct an isotopy of the annular twist.
\end{proposition}
\Uses{the hypotheses, notation and definitions stated here; \cref{push:src:C.3.1}}

\begin{proof}
Their quotient is an equivariant orbit-preserving diffeomorphism
\(h\). On \(X\setminus F\) it has a unique expression
\(h(x)=g(x)\cdot x\), with \(g\) a smooth invariant circle-valued
function. We check its extension across \(F\) directly. Choose an
equivariant tubular slice and trivialize its normal bundle locally
over \(F\). Conjugating the negative-weight complex coordinate
identifies the normal action with scalar multiplication on
\(\mathbb C^2\). Since \(h\) preserves each orbit, it fixes the
base variable \(y\) and has normal part
\[
 (f_1(y,z),f_2(y,z))=g(y,z)(z_1,z_2),\qquad z\ne0.
\]
In particular \(z_2f_1=z_1f_2\), an identity of smooth functions
also at zero. In the homogeneous Taylor expansion in the four real
normal variables, each pair \((p_1,p_2)\) therefore satisfies
\(z_2p_1=z_1p_2\). In the complexified polynomial ring with
variables \(z_1,z_2,\bar z_1,\bar z_2\), this implies
\(p_1=z_1q\) and \(p_2=z_2q\): first set \(z_1=0\) and divide
by \(z_1\), then cancel \(z_1\). This coefficient argument works
with coefficients depending smoothly on \(y\). It produces a
common formal Taylor series for \(g\).

The formal series can be realized by a smooth complex function
\(G(y,z)\). Here is the parametric Borel construction needed for
this assertion. On a relatively compact base slice, write its
homogeneous terms as \(q_n(y,z)\), choose a smooth cutoff equal to
one near zero, and replace each term by
\(\chi(z/\epsilon_n)q_n(y,z)\). Choose
\(\epsilon_n\downarrow0\) so rapidly that the term and all
mixed base/normal derivatives of total order at most \(n/2\)
are bounded by \(2^{-n}\), away from the finitely many lower
degrees. This is possible because each such derivative retains a
positive power of \(\epsilon_n\). The resulting series converges
in every smooth seminorm, and its Taylor series is the prescribed
one: the cutoff is constant near zero and higher homogeneous
terms have zero lower jets. This proves the required local
parametric extension. The remainders
\(f_j-Gz_j\) are flat in the normal variables, locally uniformly
with all derivatives in \(y\). For \(z\ne0\),
\[
 g-G=\frac{\bar z_1(f_1-Gz_1)+\bar z_2(f_2-Gz_2)}{|z|^2}.
\]
A smooth function flat along \(z=0\), divided by \(|z|^2\),
extends smoothly and flatly by zero: the usual quotient derivative
formula and Taylor bounds of arbitrarily high order prove this for
every derivative. Thus \(g\) extends smoothly across the slice.
Its absolute value is one by continuity. Equivariance implies
invariance on the free part, hence on the slice by continuity.
The local extensions agree because they agree on the dense free
stratum. We obtain a global smooth invariant \(g:X\to S^1\).

Since \(X\) is simply connected, the covering
\(\mathbb R\to S^1\), \(a\mapsto e^{2\pi ia}\), gives a smooth
lift \(a\). The integer-valued difference
\(a(\lambda x)-a(x)\) is zero at \(\lambda=1\) and continuous
in the connected circle, so \(a\) is invariant. The maps
\(x\mapsto e^{2\pi isa(x)}\cdot x\), \(0\leq s\leq1\),
are diffeomorphisms with the explicit inverse obtained by changing
\(s\) to \(-s\). They give the required isotopy.

For the second assertion the free quotient is a principal circle
bundle. A connection lifts the compactly supported time-dependent
vector field of the orbit isotopy to an invariant vector field.
Its flow exists for the full finite time interval, is equivariant,
and extends by the identity near \(F\). The endpoint differs from
any other smooth equivariant lift of the same orbit map by the
orbit-preserving diffeomorphism just treated.

The general orbit-preserving statement is Haefliger and Salem
(1991, Theorem~3.1). We used the direct slice argument because the
non-equivariant Lemma~3.2 in that paper, as originally printed,
requires correction: complex conjugation on \(\mathbb C\)
preserves scalar-circle orbits but its quotient
\(\bar z/z\) has no continuous extension to zero. The theorem
with equivariance is unaffected; the corrected argument is also
explained in Karshon's research slides, which reproduce the original
statement. No conclusion about smooth maps of the topological orbit
space near \(F\) is assumed here.
\nocite{HaefligerSalem1991D,KarshonOrbitSlidesD}
\end{proof}

